% Generated draft: requires mathematical and source review.
\documentclass[11pt]{article}
\usepackage[margin=25mm]{geometry}
\usepackage{fontspec}
\setmainfont{Latin Modern Roman}
\usepackage{amsmath,amssymb,mathtools,unicode-math}
\setmathfont{Latin Modern Math}
\usepackage{enumitem,needspace,xurl,hyperref}
\setlength{\parindent}{0pt}
\setlength{\parskip}{5pt}
\setlength{\emergencystretch}{3em}
\newcommand{\N}{\mathbb{N}}
\newcommand{\Z}{\mathbb{Z}}
\newcommand{\Q}{\mathbb{Q}}
\newcommand{\R}{\mathbb{R}}
\newcommand{\C}{\mathbb{C}}
\newcommand{\F}{\mathbb{F}}
\newcommand{\A}{\mathbb{A}}
\newcommand{\PP}{\mathbb{P}}
\newcommand{\E}{\mathbb{E}}
\newcommand{\eps}{\varepsilon}
\newcommand{\dd}{\,\mathrm{d}}
\begin{document}
\section*{UnsolvedMath ID: 30000360}
\textbf{Source problem code:} catalogue entry
\subsection{Definitions and mathematical statement}
Let $X$ and $Y$ be smooth projective varieties over the common ground field intended by the source, birational to one another and with trivial canonical class in the intended sense. Birational means that some dense open subvarieties of $X$ and $Y$ are isomorphic. Write $D^b(\operatorname{Coh}X)$ for the bounded derived category of coherent sheaves, obtained from bounded complexes of coherent sheaves by inverting quasi-isomorphisms.

The conjecture attributed to Kawamata is
\[
 X\text{ birational to }Y,\quad K_X=0,\quad K_Y=0
 \quad\Longrightarrow\quad
 D^b(\operatorname{Coh}X)\simeq D^b(\operatorname{Coh}Y),
\]
where $K_X=0$ and $K_Y=0$ abbreviate the source's unresolved canonical-triviality convention. The equivalence is understood as an exact equivalence of the derived categories, and as linear over the common ground field under the usual convention. Existence of such an equivalence is called $D$-equivalence.

A standard strong interpretation of trivial canonical class is $\omega_X\cong\mathcal O_X$ and $\omega_Y\cong\mathcal O_Y$. The extraction does not distinguish this from torsion canonical class, numerical triviality, or vanishing of a cohomological first Chern class. These conditions must not be silently identified. It also omits the ground field and its characteristic; a complex or characteristic-zero scope cannot be inferred as an explicit hypothesis solely from the short wording.

The statement asks for the birational-to-derived-equivalent direction, for all varieties in its intended class and without a dimension bound supplied here. It does not assert the converse, does not require every chosen birational correspondence to produce a specified equivalence, and does not add vanishing of intermediate structure-sheaf cohomology or a narrower Calabi--Yau definition. The supplied assessment mentions nonbirational derived-equivalent Hilbert schemes precisely as examples that do not refute this direction.
\subsection{Short English statement}
The conjecture asks whether birational smooth projective varieties with the intended trivial canonical class necessarily have equivalent bounded derived categories of coherent sheaves. The ground field and exact meaning of canonical triviality need confirmation; the converse is not asserted.
\subsection{Sources}
\begin{itemize}
\item Local UnsolvedMath registry, permanent ID 30000360, code catalogue entry, ``Kawamata Derived-Equivalence Conjecture.'' The original statement and assessment attribute the conjecture to Kawamata; the proposer field is empty. Proposed year 2005.
\item \emph{Mini-Workshop: Heterotic Strings, Derived Categories, and Stacks} (2006), Oberwolfach Reports, pp.\ 3019--3060, DOI 10.4171/owr/2005/53; \url{https://doi.org/10.4171/owr/2005/53}. The displayed report year 2006 differs from the DOI and proposed year 2005.
\item Supplied reference labeled ``derived-equivalent non-birational Hilbert schemes'': \url{https://doi.org/10.1007/s00209-019-02281-1}.
\item Statement status corrected\_verified; the supplied verification says a source label was removed and equivalence notation and punctuation normalized without changing the conjecture.
\item Literature triage checked 2026-08-21: open. It reports important cases such as threefold flops and hyperk\"ahler varieties, says no general proof or counterexample was found in the supplied search, and notes that nonbirational derived-equivalent Hilbert schemes concern the other implication. These dated claims and their exact hypotheses require independent verification. Created 2026-08-25; updated 2026-09-16.
\item Registry curation by Ulam AI / UnsolvedMath Contributors, CC BY 4.0; underlying report and literature provenance remain as supplied. Unreviewed statement-only draft.
\end{itemize}
\end{document}
